% !TeX root = ../../Book3.tex
\chapter{基本群与解析延拓}
\label{b3:ch:09}
\BookChapterNumber{b3:ch:09}
\begin{chapterabstract}
本章把解析延拓从形式换公式变为沿路径传递函数芽。有限链稳定性给同伦不变性，单值性定理在所有路径可延拓与单连通两项条件下生成全局函数。对数、幂函数、反函数和自然边界的例子说明路径依赖与延拓障碍如何分别出现。
\end{chapterabstract}

\begin{learninggoals}
\begin{itemize}
\item\label{b3:item:09-goal-1}以端点固定同伦和拼接定义基本群，计算去心平面的基本群。
\item\label{b3:item:09-goal-2}区分函数值、函数芽和沿路径的有限代表链。
\item\label{b3:item:09-goal-3}完整证明延拓唯一性、同伦不变性与单值性定理。
\item\label{b3:item:09-goal-4}解释分支层数、临界值和自然边界，避免将局部逆误认为全局覆盖。
\end{itemize}
\end{learninggoals}

一个幂级数在收敛圆外可能仍描述同一个解析对象，但必须改变局部表示。沿不同路径绕过奇点，回到原点时有时得到不同表示；这种差别首先是函数芽的差别，然后才表现为数值差别。

\ProseParagraphBreak
前提是第1章恒等定理与局部逆、以及第5章的小盘取根。路径的同伦概念在本章从定义开始建立；证明不把任何习题中的拓扑事实作为唯一前置。

% !TeX root = ../../Book3.tex
\section{曲线与同伦}
\label{b3:sec:0901}

解析延拓沿路径传递局部函数。两条路径是否可以连续变形、变形时是否经过奇点，将决定它们传递到终点的函数芽是否相同。本节把这些拓扑数据固定下来。

\subsection{路径}
\label{b3:subsec:090101}

路径指连续映射 \(\gamma:[0,1]\rightarrow U\)。起终点为 \(\gamma(0),\gamma(1)\)，反向路径记 \(\gamma^{-1}(t)=\gamma(1-t)\)。若 \(\gamma\) 的终点是 \(\eta\) 的起点，拼接 \(\gamma*\eta\) 先走 \(\gamma\) 再走 \(\eta\)，分别在两半区间线性重参数化。

\ProseParagraphBreak
沿连续路径作解析延拓不需要曲线可求长，因为传递的是重叠邻域上的函数芽，而不是沿曲线积分。需要积分时，平面区域中的连续路径可以在固定端点下用有限折线路径逼近；把紧路径像覆盖于有限小圆盘中，逐段以线段替换便给出这种逼近及圆盘内同伦。

\subsection{同伦}
\label{b3:subsec:090102}

两个同起终点的路径端点固定同伦，指存在连续 \(H:[0,1]^2\rightarrow U\)，使 \(H(0,t)=\gamma(t)\)、\(H(1,t)=\eta(t)\)，且 \(H(s,0),H(s,1)\) 分别是固定端点。反射参数、拼接参数与恒定变形证明同伦是等价关系。

\ProseParagraphBreak
自由同伦允许闭路基点移动；它与端点固定同伦不同。基本群使用后一种。沿着移动基点的路径传递闭路类时，会出现共轭，而不是原样保留每个群元素。对整数绕数这项差别被交换性遮住，对两个以上孔洞则不能忽略。

\begin{proposition}\label{b3:prop:09-path-homotopies}
拼接在同伦类上满足结合律，常路径为单位，\(\gamma*\gamma^{-1}\) 同伦于常路径。
\end{proposition}
\begin{proof}
三个路径的不同括号仅改变各段占用的参数区间，连续移动分割点即给结合律同伦。常路径占用的参数区间可连续缩到长度0，给单位关系。往返路径可先只走 \(\gamma\) 到参数 \(1-s\)，再沿原路返回；当 \(s\) 从0变到1时，往返部分连续缩成起点，给逆元关系。
\end{proof}

\subsection{基本群}
\label{b3:subsec:090103}

\begin{definition}\label{b3:def:09-fundamental-group}
以 \(a\in U\) 为基点的闭路径端点固定同伦类组成\term[jiben-qun]{基本群}[fundamental group] \(\pi_1(U,a)\)，乘法由路径拼接给出。
\end{definition}
前一命题保证这确为群。区域路径连通，换基点的连接路径给群同构，但同构依赖所选连接路径，两个选择可能相差一个内自同构。单连通等价于任意基点的基本群平凡。

\ProseParagraphBreak
凸区域中的闭路由直线收缩 \(H(s,t)=(1-s)\gamma(t)+sa\)，故基本群平凡。若区域仅星形，同样使用星形中心，再用基点连接路径换基点。开集连通保证路径连通，可以用“能由有限内部线段到达的点集同时开且闭”证明，不需要假定区域本身凸。

\ProseParagraphBreak
换基点时，群同构并非完全没有选择。若 \(\eta\) 从 \(a\) 到 \(b\)，把基于 \(b\) 的闭路 \(\gamma\) 送到 \(\eta*\gamma*\eta^{-1}\)，得到基于 \(a\) 的类。换另一条 \(\eta\)，所得同构相差一个内自同构。因而只讨论群的抽象同构类型时可不记基点，讨论具体闭路对芽的作用时却须记录基点及连接路径。两处都写“绕某个孔一圈”，若未说明如何从基点走到孔附近，可能代表不同的共轭类。

\ProseParagraphBreak
拼接路径的参数变化不会影响其同伦类。把每段路径分配不同的时间长度，线性插值各分段时刻即可构造同伦；插入常路径也可用同样办法消去。逆路径使 \(\gamma*\gamma^{-1}\) 可收缩，具体同伦让去程只走到 \(\gamma(1-s)\) 后立即沿原路返回。群公理因此来自可写出的路径变形，而不是依赖参数拼接严格满足结合律。

\subsection{绕行}
\label{b3:subsec:090104}

\begin{proposition}\label{b3:prop:09-punctured-fundamental}
\(\pi_1(\mathbb C^*,1)\) 由绕数同构于 \(\mathbb Z\)。
\end{proposition}
\begin{proof}
径向变形 \(z\mapsto ((1-s)|z|+s)z/|z|\) 将去心平面收缩到单位圆并固定单位圆。圆周路径有连续辐角提升 \(\gamma(t)=e^{\mathrm{i}\theta(t)}\)，可沿有限小弧依次选择，指定 \(\theta(0)=0\) 后唯一。闭路给 \(\theta(1)=2\pi n\)。端点固定同伦中的终点整数连续不变；反之具有相同 \(n\) 的两个辐角提升可作线性插值，指数后给端点固定同伦。因此类由整数唯一决定，拼接使整数相加。
\end{proof}
绕两个不同点的区域不再由一个整数分类。绕数只记录每个点的总绕行次数，而不记录绕行顺序；两个孔的基本群中交换子可以非平凡，却对每个外点绕数为零。第10章将区分这类同伦信息与较弱的同调信息。

% !TeX root = ../../Book3.tex
\section{函数芽与解析延拓}
\label{b3:sec:0902}

一个全纯函数的全部导数只描述当前点附近的行为。函数芽把这组局部信息作为一个对象，重叠邻域的恒等关系则提供向外传递的规则。

\subsection{函数芽}
\label{b3:subsec:090201}

\begin{definition}\label{b3:def:09-germ}
在 \(a\) 的邻域全纯的函数按“在更小邻域相同”取等价类，称为在 \(a\) 的\term[hanshu-ya]{函数芽}[germ of a function]，记 \([f]_a\)。
\end{definition}
等价由局部性容易核对。两全纯芽相同，当且仅当在 \(a\) 的 Taylor 系数逐项相同，因为两级数在共同小盘收敛。仅值相同不够：\(1\) 与 \(e^z\) 在0取同值，芽却不同。仅有限多个导数相同也不够：\(0\) 与 \(z^{N+1}\) 的前 \(N\) 阶导数在0均相同。

\ProseParagraphBreak
芽可相加相乘并求导；非零芽有亚纯倒数，只有代表在中心不为零时才有全纯倒数。把函数的某个具体定义域忘掉并不意味着忘掉中心位置，两个不同中心的芽不能直接判作同一个芽。

\ProseParagraphBreak
相同中心处仅有相同值不能确定同一个芽。\(f(z)=z\) 与 \(g(z)=z+z^2\) 在0取相同值、相同一阶导数，但芽不同。所有阶导数相同则由 Taylor 展开保证芽相同。对解析延拓来说，匹配条件是重叠开集上的函数相等，等价于在该重叠连通分支上有一个聚点集的值相等；不能用有限多个插值点取代。

\ProseParagraphBreak
延拓可以超出最初 Taylor 圆盘的收敛半径。\(1/(1+z^2)\) 在0的 Taylor 级数半径为1，但有理表达式在 \(\mathbb C\setminus\{\mathrm{i},-\mathrm{i}\}\) 上全纯，可以沿许多越过单位圆的路径延拓。圆盘收敛半径由离中心最近的奇点限制，不意味着整个圆周都是自然边界。后一节的稀疏幂级数则真有自然边界，两类情况必须用沿路径的局部代表来区分。

\subsection{沿路径延拓}
\label{b3:subsec:090202}

初始芽 \([f]_a\) 沿 \(\gamma\) 的\term[jiexi-yantuo]{解析延拓}[analytic continuation]，是沿路径给一列局部全纯代表，使相邻代表在路径连接处的一个邻域相同，首芽为 \([f]_a\)。紧性允许使用有限子区间和有限小圆盘描述；在每段内，一个代表同时给该段所有位置的芽。

\ProseParagraphBreak
例如初始 \(\log z\) 芽在1取0，可以沿任何去心路径延拓。每个小盘内用 \(1/z\) 的原函数作对数，再以重叠处值匹配常数。绕0一周后，末芽变成初芽加 \(2\pi\mathrm{i}\)。初始 \(1/(1-z)\) 在0的小盘表示式也可沿避开1的路径继续到圆盘外；初始 Taylor 圆的边界不是每一点都真的奇异。

\ProseParagraphBreak
有限链并不要求一个 Taylor 级数走完路径。先在初始盘内取代表，再在相交盘内取与它重叠相等的代表，直到端点。每一次换中心都重算局部收敛盘；连续路径的紧性保证在一项已经存在的延拓上，选出的代表和衔接区可以有限覆盖。这个有限性用于稳定性证明，但不能反过来凭任意盘覆盖宣布代表存在。盘覆盖只给拓扑位置，真正的解析代表仍需函数表达式、方程或其他存在论证。

\ProseParagraphBreak
在有理函数例子中，代表由同一个有理式给出，所以路径只要避开极点就总可延拓。对于根函数，路径必须避开使局部根分支消失的零点，或者允许改用分歧曲面；对于稀疏幂级数，外侧代表可能根本不存在。因此判断“能否沿每条路径延拓”应当先分析阻碍集合，而不是先使用单值性定理来掩盖这项存在条件。

\subsection{唯一性}
\label{b3:subsec:090203}

\begin{lemma}\label{b3:lem:09-continuation-unique}
沿同一条路径的解析延拓若存在，给定初始芽后唯一。
\end{lemma}
\begin{proof}
两延拓在起点附近相同。令 \(E\) 为沿参数已经相同的位置集合；局部代表在邻域相同使 \(E\) 开。若相同位置趋于某个参数，取该参数附近两套代表的小共同圆盘，相同位置的芽使两代表在其中有开集相同，恒等定理使中心芽也相同，所以 \(E\) 闭。区间连通且 \(E\) 非空，故是全部区间。
\end{proof}
\begin{lemma}[有限链稳定性]\label{b3:lem:09-continuation-stability}
若初始芽可沿 \(\gamma\) 延拓，则充分接近 \(\gamma\) 的路径也可使用同一组有限局部代表延拓；若两路径起终点相同，末芽相同。
\end{lemma}
\begin{proof}
把路径分成有限小段，使每段的像紧含于一个代表圆盘，且相邻接点紧含于两个代表相同的共同小邻域。有限个正的距离下界给一个共同 \(\varepsilon>0\)。与原路径逐点距离小于 \(\varepsilon\) 的新路径仍落在相同段盘中，接点也在相同共同邻域，所以可逐段用同一代表。终点固定时最终使用同一末代表，故芽相同。
\end{proof}
稳定性只对已有延拓的固定紧路径给一个邻域。它不说明越过某个奇点的路径也可继续，或任意长路径都有统一延拓半径。

\subsection{最大解析延拓}
\label{b3:subsec:090204}

从一个初始芽出发，把沿允许路径能到达的全部芽收集起来。以局部代表 \(f\) 的图式 \(z\mapsto[f]_z\) 作为坐标片，得到一个投影到基区域的局部双全纯空间；函数值映射在这些坐标片上就是 \(f(z)\)。这称为该芽的最大解析延拓空间，最大性指任何其他沿同一区域的延拓都通过芽唯一地映入它。

\ProseParagraphBreak
不同基点可用不交基圆盘分开；同一基点的不同芽也可分开，因为若两个小盘代表在任一共同开子集相同，恒等定理迫使中心芽相同。因此芽空间是 Hausdorff 的。通过有理中心和有理半径的有限延拓链可选出可数坐标基，所以它是一张通常意义的 Riemann 曲面。其投影局部单射；要把代数函数的分歧点补入，则需另外完成曲面，投影将不再是局部单射。

\ProseParagraphBreak
对 \(\sqrt z\)，在去心平面上芽空间是二层覆盖，可用 \(w\in\mathbb C^*\)、\(z=w^2\) 表示。补入 \(w=0\) 后得到整张平面，但 \(w\mapsto w^2\) 在0分歧。补点后的对象有用，却不应与“基变量 \(z\) 的全纯芽延拓”混为一谈。

\ProseParagraphBreak
最大延拓依赖所允许的对象。若只允许全纯芽，\(1/z\) 的芽曲面不能在0上方加一个全纯点；若允许球面值的亚纯芽，就可以在0补一个取无穷远值的点。类似地，\(\sqrt z\) 在 \(\mathbb C^*\) 上的全纯芽空间投影无分歧，而用 \(z=t^2\)、\(F=t\) 补上 \(t=0\) 后，得到延伸到平面底层的分歧曲面。底层区域扩大、允许极点或允许分歧，是三个不同操作，不能都含糊称为“补一个点”。

\ProseParagraphBreak
芽空间的局部图由一个代表 \(f\) 在小盘中的全部芽 \(\{f_z\}\) 构成，投影在这一图上是一一对应的小盘坐标。若两个不同芽在同一中心，选共同小盘上的代表；恒等定理使它们要么处处同芽，要么在该小盘每个中心都不同。因此两张小盘图可分开，给 Hausdorff 性。可达性还使整个空间连通：沿一条延拓路径把有限链逐段拼成芽空间中的路径即可。局部双全纯投影只有小图信息，并不自动满足覆盖映射要求的完整均匀层分解；后面将再次区分。

% !TeX root = ../../Book3.tex
\section{单值性}
\label{b3:sec:0903}

单值性定理不声称每个芽都能延拓。它要求沿指定区域内所有路径可延拓，然后以同伦控制路径依赖。在单连通区域中，这种依赖完全消失。

\subsection{同伦不变性}
\label{b3:subsec:090301}

\begin{theorem}\label{b3:thm:09-homotopy-continuation}
设初始芽可沿区域 \(U\) 内从 \(a\) 出发的每条路径延拓。端点固定同伦的两条路径给相同末芽。
\end{theorem}
\begin{proof}
设 \(H(s,t)\) 为同伦。对每个 \(s_0\)，路径 \(H(s_0,\cdot)\) 有延拓；有限链稳定性和 \(H\) 的一致连续性给 \(s_0\) 的一个邻域，其中所有路径使用相同有限代表，末芽相同。因此末芽随 \(s\) 局部恒定。参数区间连通，故从0到1始终相同。
\end{proof}
条件“每条路径可延拓”在证明中用于每一个 \(s_0\)。若同伦扫过一个分支奇点，某个中间路径的延拓不存在，局部恒定论证就断开；不能只核对最初与最后两条路径都能延拓。

\ProseParagraphBreak
有限链稳定性需要端点固定，否则“末芽相同”没有共同中心可以比较。若只把端点轻微移动，正确结论是末芽由同一个终端代表在各自中心给出。端点固定同伦恰使这些终端代表落在同一个芽集合里，局部恒定才有意义。自由同伦的闭路若基点移动，monodromy 作用要先沿基点移动路径作共轭识别；这与基本群换基点的规则一致。

\ProseParagraphBreak
对数给一个可以手工核验的同伦过程。路径沿小盘移动时，每个局部对数的差为 \(2\pi\mathrm{i}\) 整数，连续变形不能突然改变该整数，所以绕数固定。只要同伦避开0，末分支始终一致；若允许同伦穿过0，一个中间路径没有对数芽，整数便可以在这一失效处改变。数学陈述中应把“在 \(U\) 中同伦”写全，不能只说两个图形可以在整个平面拉成一样。

\subsection{单值性定理}
\label{b3:subsec:090302}

\term[dan-zhi-xing-dingli]{单值性定理}[monodromy theorem]由同伦不变性消除单连通域上的路径依赖。

\begin{theorem}[单值性定理]\label{b3:thm:09-monodromy}
设 \(U\) 单连通，初始芽可沿 \(U\) 内每条从中心出发的路径延拓。则所有延拓芽由 \(U\) 上同一个全纯函数给出。
\end{theorem}
\begin{proof}
给 \(z\in U\)，选一条从 \(a\) 到 \(z\) 的路径，并以末芽的值定义 \(F(z)\)。单连通性使任意两条这样的路径端点固定同伦，前一定理保证定义不依赖选择。固定 \(z\) 后，末代表在小盘内全纯；把到 \(z\) 的路径再接一条小盘内线段，就给附近点的值，故 \(F\) 在该盘等于这个代表。因此 \(F\) 全纯且具有所有所需芽。
\end{proof}
若区域不单连通，延拓仍由路径同伦类控制，闭路类作用在可达芽上，称为单值群的作用或 monodromy 作用。它可以平凡，例如单值有理函数，也可以非平凡；基本群非平凡不强迫每个函数都多值。

\begin{proposition}[线性微分方程的延拓]\label{b3:prop:09-linear-ode}
设 \(a,b\) 在区域 \(U\) 全纯，\(z_0\in U\)、\(c\in\mathbb C\)。初值问题
\[
 y'+a(z)y=b(z),\qquad y(z_0)=c
\]
在 \(z_0\) 有唯一全纯解芽，并可沿 \(U\) 内每条路径延拓。若 \(U\) 单连通，这些芽由同一个全纯函数给出。
\end{proposition}
\begin{proof}
小盘内选 \(A'=a\)、\(A(z_0)=0\)，再选 \(B'=e^A b\)、\(B(z_0)=0\)。则
\[
 y=e^{-A}(c+B)
\]
直接满足方程与初值。两个解之差乘 \(e^A\) 后导数为0，因此唯一。给定路径的紧像可用有限圆盘链覆盖；在每段上取局部积分因子，以前一段重叠处的值为新初值继续。局部表达式在整个圆盘全纯，有限链不会出现有限步的解爆破。唯一性使重叠代表相等，故得到沿全路径的延拓。单连通时应用单值性定理即得全局函数。
\end{proof}
这里线性性确实保证每个圆盘上任意初值都可解到整个盘。非线性方程未必如此。例如 \(y'=y^2\)、\(y(0)=1\) 的解为 \(1/(1-z)\)，即使方程系数在整个平面全纯，解也在1出现极点。因此“系数全纯”本身不够保证所有解沿全部路径以全纯芽延拓。

\ProseParagraphBreak
在非单连通域，线性方程的闭路作用是仿射变换。沿路径续接积分因子后，\(A\) 在闭路返回时增加周期 \(\alpha=\int a\,\mathrm{d}z\)。于是齐次解 \(e^{-A}\) 乘 \(e^{-\alpha}\)；非齐次项另给一个与初值无关的常数，末初值形如 \(c\mapsto e^{-\alpha}c+\beta\)。两条闭路的作用按仿射变换复合，它们可能不交换。对于 \(y'=1/z\)，作用是平移 \(c\mapsto c+2\pi\mathrm{i}n\)；对于 \(y'+(\nu/z)y=0\)，作用是乘 \(e^{-2\pi\mathrm{i}\nu n}\)。前者没有单值解，后者在 \(\nu\in\mathbb Z\) 时所有解单值，在 \(\nu\notin\mathbb Z\) 时零解仍单值。单值性可以取决于初值，而不只是方程。

\subsection{多值函数}
\label{b3:subsec:090303}

“多值函数”可以精确地理解为一张芽曲面上的单值全纯函数，同时投影到底层区域。沿闭路延拓改变的是位于同一底层点上的芽，不是一个函数在同一空间点取多个值。对 \(z^{1/m}\)，绕0一周乘 \(e^{2\pi\mathrm{i}/m}\)，故有 \(m\) 个芽；对 \(\log z\)，绕行产生无限多个相差 \(2\pi\mathrm{i}n\) 的芽。

\ProseParagraphBreak
有些函数确有自然边界。例如 \(F(z)=\sum_{n\ge0}z^{2^n}\) 在单位盘全纯，单位圆上每个 \(2^m\) 次单位根 \(\zeta\) 都是不可延拓点。当 \(r\uparrow1\)，从 \(n\ge m\) 起有 \((r\zeta)^{2^n}=r^{2^n}>0\)，其和趋于无穷：给定任意有限项数，先令这些项都接近1，再令项数增加。前面有限项保持有界，所以 \(F(r\zeta)\) 无界。这些根在圆周稠密；若能跨某一圆周点延拓到小盘，该盘与圆周的一段开弧内就不应有这样的奇点，矛盾。单值性定理不能跨越没有延拓的自然边界。

\ProseParagraphBreak
把芽的全部闭路作用汇集起来，可以看到函数与基本群之间的信息损失。某些不同闭路作用相同：平方根只记录绕数的奇偶，对数只记录一个整数，而双孔域的基本群还保留非交换次序。一个群元素对初始芽作用平凡，也未必对全部可达芽作用平凡；要定义整个芽纤维上的置换作用，必须追踪每个芽的继续延拓。后面的覆盖 monodromy 则以完整纤维为对象，每条闭路给该纤维的一个双射。

\ProseParagraphBreak
自然边界例子的结论也强于“Taylor 级数不收敛到外面”。在任意一个圆周点延拓，代表的小盘会碰到一个开圆弧；稠密的不可延拓点落在该弧内，限制回原盘会与原级数相等，因而也会延拓那个已知奇点，矛盾。这里恒等定理负责把新代表与原函数接起来，稠密性负责把局部延拓变成对某个具体奇点的矛盾，两项缺一不可。

\subsection{\texorpdfstring{\(\log z\)}{log z}、\texorpdfstring{\(z^\alpha\)}{z 的 \AnalysisPdfAlpha 次幂} 与反函数}
\label{b3:subsec:090304}

对 \(\log z\)，末值变化为 \(2\pi\mathrm{i}\operatorname{Ind}(\gamma,0)\)。对 \(z^\alpha=\exp(\alpha\log z)\)，绕行 \(n\) 次乘 \(e^{2\pi\mathrm{i}\alpha n}\)。若 \(\alpha=p/q\) 为最简有理数，有 \(q\) 层；若 \(\alpha\) 非有理实数或非实复数，层数无限。非实 \(\alpha\) 的因子模也改变，不能将所有分支都画成同一个模长的旋转。

\ProseParagraphBreak
反函数芽在导数非零处由局部逆函数定理存在。指数的逆芽就是对数；多项式的逆在临界值之外可沿路径延拓，闭路可能置换根。例如 \(z^2\) 的逆绕0交换两根。对于 \(z^3-3z\)，临界点为 \(\pm1\)，临界值为 \(\mp2\)；绕这些值分别交换相遇的两个根，两个交换一般不交换。根的全部对称函数仍单值，说明多值根与单值系数可以同时存在。

\ProseParagraphBreak
寻找反函数延拓的障碍时，临界值是局部障碍，但对非多项式整函数还可能出现有限渐近值。\(e^z\) 没有临界点，却不能在0附近取得逆芽；沿实轴向负无穷的像趋于0。不能仅从导数处处非零就宣布它在整个目标平面是一张覆盖。

\ProseParagraphBreak
双平方根的例子进一步区分单个分支点与整体绕行。对 \(\sqrt{z^2-1}\)，在 \(\mathbb C\setminus\{-1,1\}\) 中绕1或绕 \(-1\) 一次都使根变号；绕两点的闭路使 \(z^2-1\) 绕0两次，因而根不变。可以在连接 \(-1\) 与1的线段外选一个全纯分支，并以 \(F(z)/z\to1\) 在无穷远附近确定它。另一个分支为 \(-F\)。两根的平方和系数始终单值，不代表每个根都单值。

\ProseParagraphBreak
还要区分函数值回来和芽回来。一个闭路延拓后，末值偶然等于起值，末导数却可能不同；判定单值应比较末芽。对于多项式逆芽，起点取一个正则值，不同根本身不同，所以值足以区分这些逆芽；对于一般延拓问题没有这种保证。若在一个有聚点的小邻域中全部值回来，恒等定理才把值的比较提升为芽的比较。

\ProseParagraphBreak
整函数的有限渐近值还给逆芽的非紧致障碍。\(\int_0^z e^{-t^2}\,\mathrm{d}t\) 的导数处处非零，但沿正实轴趋于有限极限。某个逆分支沿像曲线向这个极限逼近时，源点逃向无穷，不能在该极限处接成有限源点的局部逆。局部逆函数定理没有失效；失效的是把局部逆图串联成整个目标上的完整覆盖所需的全局控制。多项式没有有限源逃逸，因为其适当性由最高次项增长保证，下一章覆盖论将利用这点。

\ProseParagraphBreak
割线选择可以改变一个分支的定义域，却不改变全部延拓的芽曲面。\(\log z\) 可在负实轴外定义主分支，也可沿任一从0通向无穷远的合适简单割线定义一个分支；这些单值域各自只选择同一无限层曲面上的一块。割线两侧的值相差 \(2\pi\mathrm{i}\)，是所选局部截面在重叠处的差，并非函数在曲面上不连续。画图时割线是为选分支而引入，未必是原函数本身的一条自然奇性边界。

\ProseParagraphBreak
对于 \(z^\alpha\)，若 \(\alpha=p/q\) 约分后有 \(q\) 层，其非零芽曲面可用参数 \(t\) 表示 \(z=t^q,F=t^p\)。若 \(p<0\)，\(t=0\) 处 \(F\) 为极点；若 \(p>0\)，为零点。底层映射在那里分歧，取决于 \(q\)，而函数本身的零极点阶取决于 \(p\)。分支重数与函数值阶数不是同一个整数，这在代数曲线的局部坐标中尤其重要。


\begin{chaptersummary}
解析延拓存在性、同一路径唯一性和路径独立性是三件事。有限链稳定性把已有延拓与路径小扰动连接起来，同伦与单连通性再消除全局依赖。最大芽空间上的函数单值，底层投影可能多层；补入分歧点需另作曲面完成。
\end{chaptersummary}

\BookExercises

% 本章题组自拟；以正文已建立的工具设计计算、反例、构造与综合训练。
\begin{exercise}\label{b3:ex:09-1}
计算去心平面中 \(\gamma(t)=2e^{6\pi\mathrm{i}t}\) 的绕数，求对数芽沿它绕行的变化。
\end{exercise}
\begin{solution}
绕数为3，因此对数末芽比初芽增加 \(6\pi\mathrm{i}\)。径向基点连接不改变整数；它只把基点由1换到2。
\end{solution}

\begin{exercise}\label{b3:ex:09-2}
证明两个在 \(a\) 的全纯芽若全部导数相同，则芽相同；有限阶导数为何不够？
\end{exercise}
\begin{solution}
它们 Taylor 系数逐项相同，两个级数在共同小盘收敛，故代表在该盘相同。有限阶时 \(0\) 与 \((z-a)^{N+1}\) 给反例。
\end{solution}

\begin{exercise}\label{b3:ex:09-3}
求 \(z^{2/5}\)、\(z^{\sqrt2}\)、\(z^{\mathrm{i}}\) 在去心平面上的分支层数。
\end{exercise}
\begin{solution}
绕行因子分别为 \(e^{4\pi\mathrm{i}/5}\)、\(e^{2\pi\mathrm{i}\sqrt2}\)、\(e^{-2\pi}\)。第一因子阶数5，故5层；后二者没有正整数幂为1，故无限层。第三种的分支模长按实比例变化。
\end{solution}

\begin{exercise}\label{b3:ex:09-4}
在单连通域 \(U\) 上无零全纯函数 \(f\) 的一个局部对数芽能否形成全局对数？分别用第5章和本章说明。
\end{exercise}
\begin{solution}
第5章以 \(f'/f\) 的原函数构造全局对数。按本章，小盘中存在对数，沿任意路径用有限小盘逐次匹配常数可延拓；单值性定理给全局函数，其指数与 \(f\) 局部相同，恒等定理给全域相同。
\end{solution}

\begin{exercise}\label{b3:ex:09-5}
说明 \(\sqrt z\) 的二层芽空间与补入分歧点后的曲面各自的投影性质。
\end{exercise}
\begin{solution}
芽空间用 \(w\ne0\)、\(z=w^2\) 表示，投影导数 \(2w\ne0\)，是局部双全纯二层覆盖。补入 \(w=0\) 后函数 \(w\) 全纯，但投影在0导数为零，为分歧映射而非覆盖。
\end{solution}

\begin{exercise}\label{b3:ex:09-6}
证明 \(\sum_{n\ge0}z^{3^n}\) 的单位圆也为自然边界。
\end{exercise}
\begin{solution}
对任意 \(3^m\) 次单位根，充分晚幂项在径向趋近时为正实且其和趋于无穷。此类根在圆周稠密。若某点可跨圆周延拓，一个小邻域会包含一个这样的根，并使其附近值有界，矛盾。
\end{solution}

\begin{exercise}\label{b3:ex:09-7}
某芽能沿两条同伦路径延拓，是否足以保证两条路径末芽相同？说明还缺何条件。
\end{exercise}
\begin{solution}
仅两条路径各自可延拓不够，变形中可能穿过使延拓失效的奇点。需要整个同伦的每条路径可延拓，或更强的整个区域每条路径可延拓，才能对参数应用有限链稳定性。
\end{solution}

\begin{exercise}\label{b3:ex:09-8}
若一个芽沿所有路径延拓且每条基点闭路都回到原芽，证明即使区域不单连通仍有全局全纯代表。
\end{exercise}
\begin{solution}
两条到同一终点的路径拼成闭路。沿闭路作用平凡，再沿第二条路径延拓唯一性比较，得到两末芽相同。于是按单值性定理证明中的逐点定义和局部代表粘合构造全局函数；所需是作用平凡，而非整个基本群平凡。
\end{solution}
